\section{Ende-til-ende-testing}
\label{sec:e2e-testing}

Som et supplement til de automatiserte testene ble det 
utviklet et ende-til-ende-testskript som kjøres mot det 
deployede systemet på HDOs OpenStack-VM mot faktiske 
AntMedia- og NHN-testmiljøer. Skriptet ble generert med 
AI-assistanse basert på teamets spesifikasjon av 
testscenarier og forventet systemoppførsel, og ligger 
i prosjektets repositorium under \texttt{e2e-test.sh}.
Genereringen foregikk i flere iterasjoner der teamet 
reviderte og validerte hver fase manuelt mot reelle 
responsformater fra HDOs testmiljøer før koden ble 
committet. Assertion-logikken ble spesifikt kontrollert 
mot dokumenterte HTTP-statuskoder fra AntMedia og NHN, 
og faser som ikke samsvarte med observert API-oppførsel 
ble forkastet og regenerert.

Skriptet er organisert i ti faser som dekker 
infrastruktur, autentisering, full livssyklus for begge 
backends, feilhåndtering, rate limiting, samtidighet, 
vedvarende last og metrikk-integritet. Hver fase bruker 
assertions på HTTP-statuskoder, responsbody-innhold og 
Prometheus-metrikker.

Et eksempel fra fase 3 viser hvordan AntMedia-livssyklusen 
testes — romopprettelse, WebRTC-lenke og verifisering av 
responsfelt:

\begin{lstlisting}[language=bash, caption={E2E-test — AntMedia romopprettelse og WebRTC-lenke}]
R=$(api "POST" "/api/rooms" '{"provider":"antmedia","name":"E2E-Room"}')
assert_eq "create AM room -> 201" 201 "$(status "$R")"
ID=$(id_of "$R")
assert_has "response has id" '"id":"antmedia:' "$(body "$R")"

R=$(api "POST" "/api/rooms/$ID/webrtc-link" '{"type":"publish"}')
assert_eq "publish link -> 200" 200 "$(status "$R")"
B=$(body "$R")
assert_has "has tokenId" '"tokenId"' "$B"
assert_has "has streamId" '"streamId"' "$B"
assert_future "$ID publish" "$(numfield "expiresAt" "$B")"
\end{lstlisting}

NHN-fasen verifiserer at ikke-støttede operasjoner 
returnerer HTTP~501 som forventet:

\begin{lstlisting}[language=bash, caption={E2E-test — NHN 501-responser}]
R=$(api "POST" "/api/rooms/$ID/start")
assert_eq "NHN start -> 501" 501 "$(status "$R")"

R=$(api "POST" "/api/rooms/$ID/stop")
assert_eq "NHN stop -> 501" 501 "$(status "$R")"

R=$(api "POST" "/api/rooms/$ID/token" '{"token":"ignored"}')
assert_eq "NHN setToken -> 501" 501 "$(status "$R")"
\end{lstlisting}

Skriptet kjører i om lag 40 minutter og avsluttes med 
opprydding av alle opprettede rom og verifisering av at 
Prometheus-gauger nullstilles. Fullstendig kildekode 
ligger i prosjektets repositorium:

\begin{center}
\url{https://github.com/Gruppe-213/HDO-VideoAPI}
\end{center}

\subsection{API-testing med Postman og .http-filer}
\label{subsec:postman}

Alle endepunkter ble testet manuelt mot HDOs testmiljøer 
via \texttt{HDO-VideoAPI.http}-filen i repoet og Postman. 
Dette dekket komplette arbeidsflyter for begge backends, 
inkludert autentiseringsflyt mot Keycloak, opprettelse 
og sletting av rom, start og stopp av strøm, 
tokengenerering og WebRTC-lenker. Feilscenarier som 
ugyldige rom-ID-er, manglende tokens og ukjente 
providere ble verifisert manuelt før de ble formalisert 
som integrasjonstester.

\subsection{Prometheus og Grafana}
\label{subsec:grafana}

Prometheus- og Grafana-integrasjonen ble verifisert ved å kjøre et kontrollert 
trafikkmønster gjennom E2E-scriptet og deretter sammenligne forventede hendelser 
med verdiene i Prometheus og Grafana-dashboardene. Verifiseringen ble gjort manuelt, 
men etter en fast sjekkliste:

\begin{itemize}
    \item \texttt{hdo\_active\_rooms\_total} økte ved romopprettelse og gikk tilbake 
          til null etter sletting
    \item \texttt{http\_requests\_total} økte for de endepunktene E2E-scriptet traff
    \item \texttt{http\_request\_duration\_seconds} viste registrerte responstider 
          for begge backends
    \item rate limiting-metrikker økte når scriptet overskred definert terskelverdi
    \item health check-status for AntMedia og NHN samsvarte med tilgjengeligheten 
          observert fra \texttt{/health/ready}
\end{itemize}

\subsection{Frontend-demo}
\label{subsec:frontend-demo}

En enkel frontend ble utviklet for å demonstrere 
VideoAPI-funksjonaliteten i praksis. Gjennom 
frontend-demoen ble følgende manuelt verifisert 
ende-til-ende:

\begin{itemize}
    \item Opprettelse av videorom via begge backends
    \item Visning av guestURL og hostURL for NHN Video
    \item Start og stopp av AntMedia-strøm
    \item Henting av roominformasjon
    \item Feilmeldinger ved ugyldig input
\end{itemize}

Ende-til-ende-testingen bekreftet at begge backends 
fungerer som forventet under realistiske forhold og 
at Prometheus-metrikker oppdateres korrekt gjennom 
hele romlivssyklusen.